% =====================================================================
% ANEXOS
% =====================================================================
\chapter{Instrumentos de Recolección}
\label{ap:instrumentos}

\section{Cuestionario de Percepción (Estudiantes)}

\begin{table}[H]
\centering
\caption{Cuestionario de percepción estudiantil (extracto)}
\label{tab:cuestionario-estudiante}
\small
\begin{tabularx}{\textwidth}{@{}p{1.0cm} X c@{}}
\toprule
\rowcolor{azulClaro}
\textbf{N.º} & \textbf{Ítem} & \textbf{Escala} \\
\midrule
1 & Conozco claramente cómo solicitar un cupo en una materia. & 1--5 \\
2 & El tiempo que tarda la respuesta a mi solicitud es aceptable. & 1--5 \\
3 & Puedo comunicarme con mis docentes por un canal institucional. & 1--5 \\
4 & Recibo notificación oportuna cuando se registra una falta. & 1--5 \\
5 & La evaluación docente se aplica en un momento útil del periodo. & 1--5 \\
6 & Confío en que mi evaluación docente es anónima. & 1--5 \\
7 & Sé a dónde dirigir una queja o sugerencia. & 1--5 \\
8 & Recibo información sobre el estado de mi queja o sugerencia. & 1--5 \\
9 & El horario asignado se ajusta a mi disponibilidad real. & 1--5 \\
10 & Considero que las decisiones del departamento se basan en información. & 1--5 \\
\bottomrule
\end{tabularx}
\end{table}

\section{Guion de Entrevista Semiestructurada (Docentes)}

\begin{enumerate}[leftmargin=1.4cm, itemsep=3pt]
  \item ¿Cómo gestiona actualmente la solicitud y aprobación de cupos?
  \item ¿Qué dificultades encuentra al crear o modificar una materia?
  \item ¿Cómo se comunica con sus estudiantes fuera del aula?
  \item ¿Qué proceso sigue para registrar y notificar una falta?
  \item ¿Cómo se entera de las quejas relacionadas con su materia?
  \item ¿Qué opina sobre la aplicación actual de la evaluación docente?
  \item ¿Qué información considera que le falta para mejorar su práctica?
  \item ¿Qué esperaría de un sistema integrado de gestión departamental?
  \item ¿Qué le preocuparía de un sistema que registre su desempeño?
  \item ¿Qué funcionalidades priorizaría si solo pudiera elegir tres?
\end{enumerate}

\section{Ficha de Observación de Procesos}

\begin{table}[H]
\centering
\caption{Ficha de observación del proceso de solicitud de cupo}
\label{tab:ficha-observacion}
\small
\begin{tabularx}{\textwidth}{@{}l X@{}}
\toprule
\rowcolor{verdeClaro}
\textbf{Campo} & \textbf{Descripción} \\
\midrule
Proceso observado & Solicitud de cupo en una materia del \ac{DIC}. \\
Fecha y hora & Registro del inicio y fin de la observación. \\
Actor ejecutor & Estudiante y personal de secretaría involucrado. \\
Herramientas usadas & Correo, mensajería, hojas de cálculo, documento físico. \\
Número de pasos & Conteo de acciones manuales requeridas. \\
Interrupciones & Esperas, retrocesos y solicitudes de información adicional. \\
Tiempo total & Duración medida desde la solicitud hasta la respuesta. \\
Puntos de falla & Momentos donde se pierde información o se duplica trabajo. \\
Observaciones & Comentarios cualitativos del observador. \\
\bottomrule
\end{tabularx}
\end{table}

\section{Cuestionario de Disponibilidad y Preferencia de Horarios}

\begin{table}[H]
\centering
\caption{Formato del cuestionario de horarios}
\label{tab:cuestionario-horario}
\small
\begin{tabularx}{\textwidth}{@{}l X c@{}}
\toprule
\rowcolor{moradoClaro}
\textbf{Campo} & \textbf{Descripción} & \textbf{Tipo} \\
\midrule
Código de estudiante & Identificador anonimizado & Texto \\
Semestre & Semestre que cursa & Entero \\
Materias inscritas & Códigos de materias del periodo & Lista \\
Disponibilidad & Marca por franja horaria (L--V, 6:00--22:00) & Booleano \\
Preferencia & Nivel de preferencia por franja & 1--5 \\
Restricciones & Razones de indisponibilidad (trabajo, transporte) & Categórico \\
Modalidad & Preferencia presencial / virtual / híbrida & Categórico \\
\bottomrule
\end{tabularx}
\end{table}

\chapter{Matriz de Requerimientos y Casos de Uso}
\label{ap:trazabilidad}

\begin{table}[H]
\centering
\caption{Trazabilidad entre módulos, requerimientos y diagramas}
\label{tab:trazabilidad-completa}
\small
\begin{tabularx}{\textwidth}{@{}l l X@{}}
\toprule
\rowcolor{azulClaro}
\textbf{Módulo} & \textbf{Requerimientos} & \textbf{Artefactos de modelado} \\
\midrule
Autenticación y roles & RF-01 a RF-04 & Clase \texttt{Usuario} (Fig.~\ref{fig:clases-detalle}) \\
Cupos & RF-05 a RF-10 & UC1 (Fig.~\ref{fig:casos-uso}), clase \texttt{SolicitudCupo} \\
Materias & RF-11 a RF-15 & UC2, clase \texttt{Materia} \\
Mensajería & RF-16 a RF-20 & UC6, UC7, clase \texttt{Mensaje} \\
Faltas & RF-21 a RF-26 & UC3, clase \texttt{Falta} \\
Evaluación docente & RF-27 a RF-34 & UC10 (Fig.~\ref{fig:secuencia-evaluacion}) \\
Quejas y sugerencias & RF-35 a RF-42 & UC8, Figs.~\ref{fig:actividades-quejas} y \ref{fig:estados-queja} \\
Cuestionarios & RF-43 a RF-48 & UC5, clase \texttt{CuestionarioHorario} \\
Analítica & RF-49 a RF-56 & Clase \texttt{AnaliticaDatos} (Fig.~\ref{fig:motor-analitico}) \\
\bottomrule
\end{tabularx}
\end{table}